%%
%% This is file `sample-sigconf.tex',
%% generated with the docstrip utility.
%%
%% The original source files were:
%%
%% samples.dtx  (with options: `all,proceedings,bibtex,sigconf')
%% 
%% IMPORTANT NOTICE:
%% 
%% For the copyright see the source file.
%% 
%% Any modified versions of this file must be renamed
%% with new filenames distinct from sample-sigconf.tex.
%% 
%% For distribution of the original source see the terms
%% for copying and modification in the file samples.dtx.
%% 
%% This generated file may be distributed as long as the
%% original source files, as listed above, are part of the
%% same distribution. (The sources need not necessarily be
%% in the same archive or directory.)
%%
%%
%% Commands for TeXCount
%TC:macro \cite [option:text,text]
%TC:macro \citep [option:text,text]
%TC:macro \citet [option:text,text]
%TC:envir table 0 1
%TC:envir table* 0 1
%TC:envir tabular [ignore] word
%TC:envir displaymath 0 word
%TC:envir math 0 word
%TC:envir comment 0 0
%%
%% The first command in your LaTeX source must be the \documentclass
%% command.
%%
%% For submission and review of your manuscript please change the
%% command to \documentclass[manuscript, screen, review]{acmart}.
%%
%% When submitting camera ready or to TAPS, please change the command
%% to \documentclass[sigconf]{acmart} or whichever template is required
%% for your publication.
%%
%%
\documentclass[sigconf, natbib]{acmart}
% \usepackage{cite}
% \usepackage{amsmath,amssymb,amsfonts}
\usepackage{algorithmic}
% \usepackage{graphicx}
% \usepackage{textcomp}
% \usepackage{xcolor}
% \usepackage{lipsum}
\usepackage[ruled, linesnumbered]{algorithm2e}
% \usepackage{braket}
% \usepackage{minted}
\usepackage{multirow}
% \usepackage{colonequals}
\usepackage{amsthm}
% \usepackage{indentfirst}
\usepackage{tikz}
% \usepackage{caption}
\usepackage{subcaption}
\usepackage{etoolbox}
\usepackage{amsmath}
\usepackage{amsfonts}
\usepackage{graphicx}
\usepackage{subcaption}
\usepackage{booktabs}
\usepackage{xcolor}

\usetikzlibrary{quantikz2}

\newcounter{examplecounter}
\newcommand{\example}{%
    \stepcounter{examplecounter}%
    \textbf{Example \theexamplecounter:}\hspace{0.5em}%
}

\settopmatter{printacmref=false}

%%
%% \BibTeX command to typeset BibTeX logo in the docs
\AtBeginDocument{%
  \providecommand\BibTeX{{%
    Bib\TeX}}}

%% Rights management information.  This information is sent to you
%% when you complete the rights form.  These commands have SAMPLE
%% values in them; it is your responsibility as an author to replace
%% the commands and values with those provided to you when you
%% complete the rights form.
\setcopyright{rightsretained}
\copyrightyear{2026}
\acmYear{2026}
\acmDOI{XXXXXXX.XXXXXXX}
%% These commands are for a PROCEEDINGS abstract or paper.
\acmConference[DAC '26]{63rd ACM/IEEE Design Automation Conference}{July 26--29, 2026}{Long Beach, CA, USA}
%%
%%  Uncomment \acmBooktitle if the title of the proceedings is different
%%  from ``Proceedings of ...''!
%%
% \acmBooktitle{DAC '24: Proceedings of the 61st ACM/IEEE Design Automation Conference, July 26--29, 2026, Long Beach, CA, USA}
\acmISBN{XXX-X-XXXX-XXXX-X/26/XX}


%%
%% Submission ID.
%% Use this when submitting an article to a sponsored event. You'll
%% receive a unique submission ID from the organizers
%% of the event, and this ID should be used as the parameter to this command.
%%\acmSubmissionID{123-A56-BU3}

%%
%% For managing citations, it is recommended to use bibliography
%% files in BibTeX format.
%%
%% You can then either use BibTeX with the ACM-Reference-Format style,
%% or BibLaTeX with the acmnumeric or acmauthoryear sytles, that include
%% support for advanced citation of software artefact from the
%% biblatex-software package, also separately available on CTAN.
%%
%% Look at the sample-*-biblatex.tex files for templates showcasing
%% the biblatex styles.
%%

%%
%% The majority of ACM publications use numbered citations and
%% references.  The command \citestyle{authoryear} switches to the
%% "author year" style.
%%
%% If you are preparing content for an event
%% sponsored by ACM SIGGRAPH, you must use the "author year" style of
%% citations and references.
%% Uncommenting
%% the next command will enable that style.
%%\citestyle{acmauthoryear}


%%
%% end of the preamble, start of the body of the document source.
\begin{document}

%%
%% The "title" command has an optional parameter,
%% allowing the author to define a "short title" to be used in page headers.
\title{From Graph \texorpdfstring{$k$}{k}-Coloring to Grover's Circuits via Logic Synthesis}

%
% The "author" command and its associated commands are used to define
% the authors and their affiliations.
% Of note is the shared affiliation of the first two authors, and the
% "authornote" and "authornotemark" commands
% used to denote shared contribution to the research.
\author{Ismael Barzani}
\affiliation{%
  \institution{ECE Department - Concordia University}
  \city{Montréal}
  \state{Québec}
  \country{Canada}
}
\email{ismael.ridha@mail.concordia.ca}

\author{Sébastien Le Beux}
\affiliation{%
  \institution{ECE Department - Concordia University}
  \city{Montréal}
  \state{Québec}
  \country{Canada}
}
\email{sebastien.lebeux@concordia.ca}

% Uncomment the above for the info. Uncomment the header author info also.
% \author{Omitted-for-blind-review}
% \affiliation{%
%   \institution{Omitted-for-blind-review}
%   \city{Omitted}
%   \state{Omitted}
%   \country{Omitted}
% }
% \email{Omitted-for-blind-review}


%%
%% By default, the full list of authors will be used in the page
%% headers. Often, this list is too long, and will overlap
%% other information printed in the page headers. This command allows
%% the author to define a more concise list
%% of authors' names for this purpose.
\renewcommand{\shortauthors}{Barzani \& Le Beux}
% \renewcommand{\shortauthors}{Omitted-for-blind-review}

%%
%% The abstract is a short summary of the work to be presented in the
%% article.
\begin{abstract}
Graph $k$-coloring provides a test case for compiling Boolean predicates into quantum search circuits. We describe a logic-synthesis flow with explicit color encodings, restored oracle workspace, and a diffusion operator matched to the actual state preparation. The study asks when representation, preparation, mapping, and multi-controlled-gate decomposition change resource costs under a common semantic contract. XOR--AND--inverter graphs provide two-input logic primitives, but their use alone establishes neither a resource advantage nor novelty in automated oracle synthesis. This revision defines the comparison and its limitations; corrected comparative results, fault-tolerant costs, noisy simulations, and scaling conclusions remain pending independent evidence review.
\end{abstract}
\maketitle

\section{Introduction}
\label{sec:introduction}
A proper $k$-coloring of a graph assigns one of $k$ colors to each vertex so that adjacent vertices differ. Grover search offers a quadratic query reduction relative to exhaustive search over a chosen domain~\cite{grover}; this does not establish an advantage over the best classical coloring algorithms or account for circuit implementation costs. The practical compilation question is how to implement the predicate, preparation, and reflection with a consistent workspace and gate model.

Boolean-to-quantum compilation, including coloring oracles, is established prior work~\cite{DATE21}. Restricted color-state preparation is also established~\cite{MukherjeeList,LevinUniform}. We therefore ask: under matched correctness, domain, preprocessing, workspace, and lowering requirements, when do logic-based coloring constructions reduce resource costs, and which mechanisms explain the trade-offs? The present revision specifies the encoding and comparison methodology. Any contribution based on improved resources, scaling, or noisy success awaits the corrected evaluation. We make no priority or general superiority claim.

\section{Background and Related Work}
\label{sec:background}
\subsection{Logic representations and prior compilation flows}
\label{subsec:related-logic}
An XOR--AND--inverter graph (XAG) is a directed acyclic Boolean network with two-input XOR and AND nodes and complemented signals. It is not canonical: equivalent functions can have different networks. Multiplicative depth is the maximum number of AND nodes on an input-to-output path; AND count and multiplicative depth describe a particular network, whereas multiplicative complexity is the minimum AND count over representations of the function. These distinctions matter when relating networks to quantum costs~\cite{XAGCompilation,XAG_npj,HaenerSoekenDepth}.

Schmitt et al.~\cite{DATE21} present an automated flow from Python Boolean functions to quantum circuits, including uniform state preparation, XAG and LUT oracle synthesis, and a vertex-coloring demonstration for the Zed city problem. It is not merely a manual coloring construction. Meuli et al.~\cite{XAGCompilation,XAG_npj} study XAG compilation and trade-offs among qubits, $T$-count, and $T$-depth. H\"aner and Soeken~\cite{HaenerSoekenDepth} reduce multiplicative depth using cut enumeration, tree balancing, and exclusive sum-of-products representations. Balancing in this study follows that established line of work.

Yu et al.~\cite{YuBackend2025} connect XAG properties to the $T$ count, logical time steps, and helper qubits of back-end-compiled fault-tolerant oracles. Their back-end-aware objectives show why logical AND counts alone cannot establish implementation cost. Our intended comparison concerns coloring compilation under matched contracts; it does not reproduce their back-end model or estimate physical error-correction resources.

\subsection{Coloring and restricted state preparation}
\label{subsec:related-coloring}
Comparator-based coloring circuits are studied by Saha et al.~\cite{saha} and revised by Belletti et al.~\cite{belletti}. These provide comparison constructions rather than evidence that one representation is uniformly preferable. Campbell, Khurana, and Montanaro~\cite{CampbellColoring} analyze Grover and quantum backtracking approaches to constraint satisfaction, including graph coloring, with detailed fault-tolerant overheads. Their algorithm and cost assumptions differ from a logical-oracle comparison; their work precludes treating coloring cost analysis itself as new.

Mukherjee~\cite{MukherjeeList} uses restricted initialization and diffusion for list coloring. Levin~\cite{LevinUniform} develops uniform preparation over arbitrary domain sizes and applies it to coloring search. Shukla and Vedula~\cite{Shukla2024} give an ancilla-free preparation over consecutive basis states with logarithmic gate complexity and depth in a gate model permitting rotations and controlled Hadamards. This is not a precision-independent Clifford+$T$ cost. Preparing valid color codes and matching the reflection are prior techniques, not contributions by themselves.

Concurrent work by Sen and Mukherjee, \emph{A Quantum Framework for K Coloring of Graphs}~\cite{SenMukherjee2026}, discusses binary encoding and constraint circuits across Grover search, QAOA, and annealing. Its Section VI-F identifies its empirical evaluation as a single small simulated graph. We cite that bounded scope without inferring a comparative ranking or reproducing its reported results.

\subsection{Multi-controlled gates and workspace contracts}
\label{subsec:workspace}
We use MCX for a multiply controlled $X$ operation and MCT for its Toffoli generalization with more than two controls. A clean ancilla begins in a known state, usually $\ket{0}$; a restoring construction returns it to that state. A dirty ancilla is borrowed in an arbitrary state and must be restored coherently, including correlations with an external register. A throwaway ancilla begins clean but retains intermediate garbage until an explicitly paired inverse removes it. A conditionally clean qubit is known on a specified control branch and usable only under the associated conditions; it is not an unconditionally fresh zero qubit~\cite{KhattarGidney2025}.

Barenco et al.~\cite{Barenco1995} provide foundational multi-controlled constructions. The decomposition study includes their one-clean-ancilla construction, the borrowed-workspace construction of Iten et al.~\cite{Iten2016}, Maslov's relative-phase constructions~\cite{Maslov2016}, the no-extra-ancilla policy based on Vale et al.~\cite{Vale2024}, Gray-code synthesis, and clean restoring trees. Vale et al. treat controlled special unitaries; an MCX implementation must retain any required phase corrections. Khattar and Gidney~\cite{KhattarGidney2025} use conditionally clean ancillae in low-workspace constructions, including clean and dirty variants. Dutta et al.~\cite{toffoli_decomposition} quantify clean-space versus Toffoli-depth trade-offs. These are distinct policies with distinct assumptions, not a universal logarithmic-depth weight for every MCX.

Figure~\ref{fig:mct-decomposition} shows only a target-action tree. The actual expanded inverse must act on the same held scratch wires after a compatible intervening operation. Garbage cannot be discarded while preserving interference. Final marking and diffusion gates require restoration within their own construction. We count peak simultaneously held workspace, including the interval between a block and its inverse.

\begin{figure}
    \centering
    \begin{subfigure}[b]{0.2\textwidth}
        \centering
        \begin{quantikz} [column sep=0.1cm, wire types={q,q,q,q,q}, row sep=0.2cm, thin lines]
             \qw & \ctrl{4} & \qw \\
             \qw & \ctrl{0} & \qw \\
             \qw & \ctrl{0} & \qw \\
             \qw & \ctrl{0} & \qw \\
             \qw & \targ{} & \qw 
        \end{quantikz}
        $\equiv_{\mathrm{target}}$ \hspace{-0.5cm} \begin{quantikz} [column sep=0.1cm, wire types={q,q,q,q,q}, row sep=0.2cm, thin lines]
             \qw &\qw & \ctrl{5} & \qw & \qw & \qw\\
             \qw &\qw & \ctrl{0} & \qw & \qw & \qw\\
             \qw &\qw & \qw & \ctrl{4} & \qw & \qw\\
             \qw &\qw & \qw & \ctrl{0} & \qw & \qw\\
             \qw & \qw & \qw & \qw & \targ{} & \qw \\
             \lstick{\small $\ket{0}_{a_1}$} & \qw & \targ{} & \qw & \ctrl{0} & \qw\\
             \lstick{\small $\ket{0}_{a_2}$} & \qw & \qw & \targ{} & \ctrl{-2} & \qw
        \end{quantikz}
        \vspace{0.8cm}
        \caption{MCT with 4 controls}
        \Description{MCT Gate with 4 controls}
        \label{fig:4-MCT}
    \end{subfigure}
    \hspace{0.2cm}
    \begin{subfigure}[b]{0.2\textwidth}
        \centering
        \begin{quantikz} [column sep=0.1cm, wire types={q,q,q,q,q}, row sep=0.2cm, thin lines]
             \qw & \ctrl{5} & \qw \\
             \qw & \ctrl{0} & \qw \\
             \qw & \ctrl{0} & \qw \\
             \qw & \ctrl{0} & \qw \\
             \qw & \ctrl{0} & \qw \\
             \qw & \targ{} & \qw 
        \end{quantikz}
        $\equiv_{\mathrm{target}}$ \hspace{-0.5cm}\begin{quantikz} [column sep=0.1cm, wire types={q,q,q,q,q}, row sep=0.2cm, thin lines]
             \qw &\qw & \ctrl{6} & \qw & \qw & \qw & \qw\\
             \qw &\qw & \ctrl{0} & \qw & \qw & \qw & \qw\\
             \qw &\qw & \qw & \ctrl{5} & \qw & \qw & \qw\\
             \qw &\qw & \qw & \ctrl{0} & \qw & \qw & \qw\\
             \qw &\qw & \qw & \qw & \ctrl{4} & \qw & \qw\\
             \qw & \qw & \qw & \qw & \qw & \targ{} & \qw\\
             \lstick{\small $\ket{0}_{a_1}$} & \qw & \targ{} & \qw & \ctrl{0} & \qw & \qw\\
             \lstick{\small $\ket{0}_{a_2}$} & \qw & \qw & \targ{} & \qw & \ctrl{0} & \qw \\
             \lstick{\small $\ket{0}_{a_3}$} & \qw & \qw & \qw & \targ{} & \ctrl{-3} & \qw
        \end{quantikz}
        \caption{MCT with 5 controls}
        \Description{MCT Gate with 5 controls}
        \label{fig:5-MCT}
    \end{subfigure}
    \caption{Throwaway-ancilla trees for MCT target action. The ancillas start in $\ket{0}$ and retain garbage; an actual expanded inverse on the same scratch wires must remove it before the oracle ends. These are not restoring decompositions.}
    \Description{}
    \label{fig:mct-decomposition}
    \vspace{-0.6cm}
\end{figure}



\subsection{Gate models and static versus dynamic synthesis}
\label{subsec:cost-model}
Logical counts, common one- and two-qubit basis counts, routed native counts, and Clifford+$T$ estimates are separate metrics. Barriers and measurements do not count as unitary gates. Noisy intermediate-scale quantum (NISQ) evaluation requires a specified target, routing, and noise model; logical width alone does not establish device compatibility.

For a selected exact, static Toffoli gadget we charge seven $T/T^\dagger$ gates and three $T$ layers~\cite{Tcount}. Selected relative-phase gadgets use four $T/T^\dagger$ gates, but substitution requires proof that the full composition has the correct phases~\cite{Maslov2016}. Gate names alone do not establish that identity. Clifford+Toffoli is an intermediate representation to be lowered, not an extension of the Clifford+$T$ set. $T$-depth is computed from dependencies in the emitted circuit; it is not the sum of isolated gadget depths. Measurement-assisted logical-AND cleanup is not free uncomputation in a static unitary model.

Ross and Selinger~\cite{RossSelinger} approximate rotations by static Clifford+$T$ circuits. The synthesis algorithm's classical randomness does not make the emitted circuit a dynamic quantum procedure. Paetznick and Svore~\cite{PaetznickSvore} instead use measurement-conditioned repeat-until-success (RUS) circuits. Their expected costs require a trial-success distribution and a retry policy. RUS cannot be substituted directly for a unitary preparation and its adjoint. The separately retained $k=3$ RUS accounting is an analytical local construction, not an executed dynamic preparation or a literature claim of coloring optimality. Dynamic forward/inverse channels, branch phases, aborts, and any coherent control require separate validation.

\section{Encoding and Circuit Construction}
\label{section:hls-flow}
\subsection{Predicate and bit ordering}
\label{subsec:encoding}
Let $G=(V,E)$ be an unweighted undirected graph with an explicit ordered vertex list, including isolated vertices, and let $k\geq1$. We use $c=\max(1,\lceil\log_2 k\rceil)$ to match the maintained one-bit convention for $k=1$, and $n=c|V|$ data qubits. The variable $x_{v,i}$ is color bit $i$, least significant first, and
\begin{equation}
 b_v(x)=\sum_{i=0}^{c-1}2^i x_{v,i},\qquad
 C_{uv}(x)=\bigvee_{i=0}^{c-1}(x_{u,i}\oplus x_{v,i}).
\end{equation}
A valid code satisfies $b_v<k$; a proper coloring additionally satisfies every edge inequality. Define
\begin{align}
 f_E(x)&=\bigwedge_{(u,v)\in E}C_{uv}(x),\\
 f(x)&=\left(\bigwedge_{v\in V}[b_v(x)<k]\right)\land f_E(x).
 \label{eq:complete-predicate}
\end{align}
An empty conjunction is true. The complete predicate is used on the unrestricted binary domain. Baseline labels are normalized to consecutive indices with an explicit inverse label map, and adapters translate their most-significant-bit ordering to this least-significant-bit convention. Precolored research instances are excluded.

For $k\geq2$, each edge clause has a two-input realization with $c-1$ ANDs by De Morgan's law applied to the bit differences. A direct conjunction of nonempty edge clauses adds $|E|-1$ ANDs, giving an unoptimized upper bound $|E|c-1$ for $f_E$ when $|E|>0$. This is a construction count, not minimum multiplicative complexity, a qubit formula, or a cost for the complete predicate's domain checks. Empty edge sets and constant simplifications are handled separately. Binary comparators can use the same local conjunction structure, so this bound alone implies no advantage over comparator baselines.

\subsection{Preparation and controlled regimes}
\label{subsec:uqs-prep}
The ideal valid-color preparation on each vertex is
\begin{equation}
 A_k\ket{0^c}=\frac{1}{\sqrt{k}}\sum_{a=0}^{k-1}\ket{a},\qquad
 A=\bigotimes_{v\in V} A_k.
\end{equation}
The supported subspace contains $N=k^{|V|}$ assignments, not necessarily proper colorings. Ideal preparation, arbitrary-rotation circuit realization, and a finite-precision Clifford+$T$ approximation are distinct objects. The maintained generic preparation uses Qiskit's state-preparation interface; the specialized expression/replay path uses the Shukla--Vedula construction. They must be identified separately in evidence, rather than attributing the specialized circuit's costs to every preparation.

Three regimes isolate effects: R1 uses Hadamards with $f$; R2 uses valid-color preparation with $f$; R3 uses valid-color preparation with $f_E$. For ideal $A$, $f=f_E$ on its support. A diagonal phase oracle preserves this support and reflection about $A\ket{0}$ also preserves it, proving the redundancy of domain checks within that ideal subspace. A finite-precision preparation can leak outside it; R3 then needs an explicit error/leakage analysis. No exact ideal-amplitude certificate follows from a numerical rotation admission.

\subsection{Mappings and clean phase oracles}
\label{subsec:mappings}
The maintained methods are named by configuration: \texttt{xag} uses tweedledum XAG synthesis; \texttt{aig\_bennett} and \texttt{xag\_bennett} use the corresponding representation with Caterpillar's Bennett mapping; \texttt{klut\_bennett} uses LUT mapping followed by Bennett scheduling; and \texttt{best\_fit} uses Caterpillar's cut-based best-fit strategy. These are configurations, not representatives of every algorithm in those families. Historical defaults leave AIG Bennett unrefactored and refactor the other representations. Comparisons must request and record common preprocessing explicitly. Encoder balancing, multiplicative-complexity reduction, cut/ESOP balancing, in-place/eager scheduling, SAT pebbling, and sequential/edge-scheduled structure-aware oracles remain distinct ablations; unsupported or unadmitted strategies are not silently included.

A raw computation may retain workspace or relative phase. The required phase oracle satisfies
\begin{equation}
 O_f\ket{x}\ket{0^w}=(-1)^{f(x)}\ket{x}\ket{0^w}.
 \label{eq:phase-contract}
\end{equation}
It marks proper colorings. Compute--phase--uncompute supplies this contract when the raw computation and its actual inverse are validated. Direct kickback on $\ket{-}$ is permitted only after arbitrary-target XOR action and coherent workspace restoration are established. Basis-output agreement alone is insufficient. ``MCT-free'' applies only to an admitted oracle construction whose emitted gates have at most two controls; it does not apply to every mapping or the full Grover circuit.

\subsection{Matching reflection, order, and solution fraction}
\label{subsec:assembly}
With $S_0=2\ket{0^n}\!\bra{0^n}-I$, the reflection and iterate are
\begin{equation}
 D=A S_0 A^\dagger,\qquad Q=D O_f,\qquad
 \ket{\psi_j}=Q^j A\ket{0^n}.
 \label{eq:grover-order}
\end{equation}
Operators act right to left: initialize once with $A$, then apply $O_f$, $A^\dagger$, $S_0$, and $A$ for each iteration. In the convention where MCZ flips the all-ones state,
$S_0=-X^{\otimes n}\mathrm{MCZ}X^{\otimes n}$.
The scalar minus fixes the stated reflection convention. The diffusion MCZ generally retains $n-1$ controls and must be lowered with the same restoring policy across methods.

For uniform preparation over a domain of size $N$ with $M$ marked assignments, $p=M/N$. R1 has $N=2^n$; ideal R2/R3 have $N=k^{|V|}$. For a general preparation use its actual marked probability $p$. With $0<p<1$ and $\theta=\arcsin\sqrt{p}$,
\begin{equation}
 P_j=\sin^2((2j+1)\theta).
 \label{eq:success}
\end{equation}
Evaluate the nonnegative floor and ceiling integers around $\pi/(4\theta)-1/2$ and choose the larger success probability, breaking ties toward fewer iterations. This selects the discrete first-peak candidate, not the minimum cost per success. If $p=0$ amplification cannot help; if $p=1$ no iteration is needed. Unknown $M$ leaves costs parametric or requires a documented counting/estimation procedure. More iterations need not improve success.

For static concatenated gate counts at a fixed chosen precision,
\begin{equation}
 T_{\mathrm{search}}=T_A+j(2T_A+T_{S_0}+T_O).
 \label{eq:search-cost}
\end{equation}
Synthesize $A$ once and reuse the actual emitted inverse for $A^\dagger$. A whole-algorithm budget counts $(1+2j)$ preparation occurrences and $j$ occurrences of each other component, including preprocessing errors. The planned primary operator-error budget is $10^{-3}$, with $10^{-2}$ and $10^{-4}$ sensitivity checks. These are requested tolerances, not achieved measurements. Numerical matrix admission, norm propagation, angle/precision limits, and caller allocation guards remain those of the exact reviewed compiler bytes. The count identity is neither a $T$-depth identity nor a dynamic RUS cost formula.

\section{Evaluation Plan: Evidence Pending}
\label{sec:experiments}
\textbf{Revision status.} Corrected tables, figures, percentages, and superiority conclusions are pending saved-circuit verification and independent exact-hash claim review. Historical figures and hardware results are retained in repository history and archives, outside this revised live evaluation. No new pilot numbers replace them here. A completed queue job is not scientific admission or protocol freeze.

The full evaluation remains required. Compare the five configured methods and four comparator baselines under R1--R3 on the nine original unprecolored graphs, including only $K_3$, $K_5$, $K_{10}$, and $K_{20}$ among complete graphs. Include strongest matched baselines per graph, geometric summaries where defined, win/tie/loss counts, and all exclusions. Verify emitted circuits and recount saved artifacts before reporting resources. Failed baselines require a documented correction and review or an explicit exclusion, not an advantage assigned to another method.

Report compiler stage times and peak memory; AND/XOR/LUT counts and multiplicative depth; data, AND, constraint, MCX, and phase workspace; logical and common-basis costs; routed native two-qubit counts/depth, active physical qubits, and scheduled duration; and complete-search Clifford+$T$ counts, dependency $T$-depth, and rotation costs. Apply matched decomposition policies to comparators, domain checks, conjunctions, LUT/best-fit gates, and diffusion. The MCX study retains 3--20 controls, with explicit unsupported, timed-out, or infeasible cases. Mapping, balancing, preprocessing, and sequential/parallel scheduling need matched ablations and common qubit budgets.

Routing/noise studies require the recorded FakeTorino snapshot and software versions, five fixed seeds with identical settings across methods, and explicit routing infeasibility. Success studies retain paths with 3--8 vertices, $K_3$, $C_5$, and $P_4$ with non-power-of-two colors, sweeps through the first ideal maximum and one iteration beyond, invalid-code leakage, shot confidence intervals, random-sampling trials, and cost per successful solution. Analytical predictions, executed ideal states, and noisy shots are labeled separately. Static graph fault-tolerant costs and dynamic RUS execution/abort accounting require separate evidence; no preparation break-even is asserted.

Scaling retains controlled vertex/density/color families and the selected DIMACS scope: queen instances, myciel6/7, DSJC125.1, le450\_5a, queen8\_8, 5-FullIns\_3, and 4-Insertions\_3. Use at least five real repetitions, independent positive witnesses where available, stage/boundary timing, memory-scope labels, missing/unsupported input ledgers, and timeouts. A small compiler pilot cannot establish scalability or quantum correctness.

\subsection{Validation and admission boundaries}
\label{subsec:validation}
Exhaustive Boolean checks cover bounded small instances; larger instances need explicit satisfying witnesses, edge violations, invalid codes, and boundary cases with sampled coverage labeled. Verify equivalence after optimization, phase and restoration on superpositions, entangled-reference restoration for borrowed workspace, and actual forward/inverse cleanup for throwaway constructions. Compare lowered operators only within the declared dense-operator limit, and multi-iteration ideal execution with independent amplitude-amplification predictions.

On the current 8\,GiB host, executed ideal statevectors are capped at 18 total qubits and dense full operators at five. Guard allocations before arrays, use one numerical/compile child at a time and one compiler/BLAS/OpenMP thread, and declare per-child and cumulative time/memory budgets. Record measured RSS and its accounting scope, timeouts, failures, exact source/input/circuit hashes, commands, and unavailable measurements. Reviewed numerical static-rotation and ideal-iteration slices do not certify arbitrary ideal amplitudes, dynamic RUS, unrestricted fault tolerance, or the whole paper. Hardware execution, physical QEC estimates, precolored research, and conditional-workspace synthesis remain excluded; conditionally clean MCX decompositions remain included.

\section{Conclusion: Revision Status}
\label{sec:conclusion}
The revised construction distinguishes valid codes from proper colorings, restores oracle workspace, and matches diffusion to preparation. Established logic synthesis and restricted-state methods motivate a controlled comparative study. Comparative resource savings, noisy success, scaling behavior, and the final contribution statement remain open until their exact evidence passes independent review. This results-independent revision is not approval of the full manuscript.

\bibliographystyle{ACM-Reference-Format}
\bibliography{refs}
\end{document}
